% Principal GDZ PPN236018647, native475/printed468; HKUST476.
\clearpage\SourceLeaf{475}\ArithmeticHead{468}{L}{GEOMETRIE. BRIEFWECHSEL.}
\begingroup\small
\HistoricTitle{BEMERKUNG.}
\indent Daß bei einem ebenen Viereck, dessen Winkel kleiner als \(180^{0}\) sind, der Schnittpunkt der Diagonalen den Punkt kleinster Entfernungssumme liefert, hatte bereits \textsc{J. Fr. A. Fagnano} bewiesen: \textit{Problemata quaedam ad methodum maximorum et minimorum spectantia}, Nova acta eruditorum anni 1775, Lipsiae 1779, S. 281. Die Verallgemeinerung auf \(n\) Punkte haben \textsc{Tédenat} und \textsc{L’Huilier} behandelt, Annales de mathématiques pures et appliquées publiées par \textsc{Gergonne} \textbf{1}, 1810—11, S. 285 und 291; bei ihnen findet sich auch schon die von \textsc{Schumacher} in dem Briefe [4.] hergeleitete Gleichung (A).

\indent Die Worte: \textls[120]{hircum mulcere}, deren sich \textsc{Gauss} im ersten Absatz der Briefstelle [7.] bedient, sind eine Anspielung auf den Vers in \textsc{Vergil}s \textit{Bucolica}, Ecloga III, 91:
\begin{center}Atque idem iungat vulpes et mulgeat hircos.\end{center}

\indent Die beiden im dritten Bande des Briefwechsels zwischen \textsc{Gauss} und \textsc{Schumacher} als Nr. 528 und Nr. 529 abgedruckten Briefe (unter [7.] ist der erste Teil von Nr. 528 wiedergegeben), bilden in Wirklichkeit einen einzigen Brief, dem das Datum des 21. April 1836 zukommt; dies geht mit voller Deutlichkeit daraus hervor, daß \textsc{Schumacher} in dem Briefe Nr. 531 den \textsc{Gauss}schen Brief vom 21. April »Ihren vorletzten Brief« nennt (a. a. O. S. 53 und 54), so daß also zwischen Nr. 529 und Nr. 530 kein Brief von \textsc{Gauss} an \textsc{Schumacher} liegt, während \textsc{Gauss} auf \textsc{Schumacher}s Brief Nr. 527 vom 19. April nicht \textls[120]{vor} dem 21. April antworten konnte.

\indent Die hier abgedruckten Briefstellen sind nach den im \textsc{Gauss}archiv befindlichen Handschriften ohne jede Anderung wiedergegeben.
\par\noindent\hfill\textsc{Stäckel}.
\ClosingRule
\par\endgroup
